\subsection{\textbf{Sprint 3}}

Levando em consideração o pedido dos \ac{ages} III, decidimos que o melhor a fazer nesta \textit{sprint} seria deixá-los sem nenhuma tarefa definida, de modo que pudessem pegar a tarefa que quisessem, caso quisessem. Nesta \textit{sprint}, voltamos a adotar o método de \textit{squads}, com tarefas que se complementam.

Novamente, junto com os outros \ac{ages} IV, decidimos revisar cada história da \textit{sprint}. Estava planejado entregar toda a parte do portal de administrador, em que seria possível adicionar simulados e também questões em lote; o cronograma, em que o usuário poderia escolher os dias e horários em que gostaria de estudar; a parte de \textit{feedback} gamificado, utilizando \textit{streak} de dias e gratificando o usuário a cada questão respondida corretamente; e, por fim, foi adicionada também uma nova tarefa de simulados, para visualizar o histórico dos simulados já realizados pelo usuário.

Para evitar os problemas que tivemos nas duas \textit{sprints} anteriores, em que alguns detalhes das tarefas não estavam totalmente alinhados, fizemos uma última revisão, agora com o time todo, para que pudessem pontuar o que poderíamos alterar. Isso nos trouxe alguns \textit{feedbacks}, e nos prontificamos a fazer as mudanças, especialmente nas tarefas do meu time, que ficou responsável pelo épico de cronograma. Um dos nossos \ac{ages} III apontou que a tarefa de \textit{backend} poderia ficar muito grande caso mantivéssemos tudo junto. Com esse retorno, decidi dividi-la em duas etapas: a primeira contemplaria apenas um \ac{crud} básico para inserção e visualização do cronograma, enquanto a segunda concentraria toda a regra de negócio, conforme imaginávamos que o cronograma deveria funcionar.

No início da \textit{sprint}, entrei em contato com o \ac{ages} I que havia ficado com a tarefa de \textit{backend}, e ele me informou que estava com alguns problemas, mas que tentaria iniciá-la. Pouco depois, no entanto, avisou que precisaria se ausentar das aulas por um tempo. Isso trouxe certa preocupação quanto à entrega do cronograma, pois, sem o \textit{backend}, ela não seria possível. No mesmo dia, em aula, conversei com o \ac{ages} III do meu time para que ele assumisse a tarefa e a desenvolvesse.

Durante esta \textit{sprint}, eu também peguei uma tarefa para desenvolver, pois era uma entrega importante para as clientes e não daria tempo de concluí-la caso contrário. Fiquei responsável por refatorar a tela da trilha de tópicos, conforme havíamos combinado com as \textit{stakeholders}. O desenvolvimento da tela em si não apresentou nenhum problema, exceto por um ponto de que não nos demos conta no início: para cada tópico existente, seria necessário criar um ícone que o diferenciasse, o que acabou levando bastante tempo.

Uma semana antes da apresentação, data do nosso \textit{code freeze}, já tínhamos todas as tarefas aguardando a revisão dos \ac{ages} III, o que nos deu, como \ac{ages} IV, o indicativo que queríamos: o time havia produzido bem, e agora faltavam apenas a aprovação e correções pontuais no código. Esse processo, porém, acabou atrasando. Mesmo após passar a semana inteira cobrando, os \ac{pr} só começaram a ser aprovados com mais frequência faltando três dias para a apresentação. No fim, dei suporte aos \ac{ages} III na correção dos \textit{bugs} nos últimos dias da \textit{sprint}. Foi nesse período que notei um \textit{bug} simples em um \textit{endpoint} de administrador: o \textit{backend} estava esperando receber o \ac{uuid} de disciplinas e tópicos, o que não fazia sentido com a regra de negócio que havíamos montado. Como isso já havia sido \textit{mergeado} na \textit{develop}, criei uma tarefa de \textit{bug} para que esse comportamento fosse corrigido na \textit{sprint} seguinte.

A apresentação para as \textit{stakeholders} ocorreu bem, porém, ao final, identificamos um \textit{bug} na tela de progresso, onde fica a parte de gamificação: ela estava considerando apenas as questões que o usuário havia respondido, e não todas as questões do simulado, como era o esperado. Infelizmente, foi uma das \textit{stakeholders} que percebeu. Anotamos como \textit{bug} para ser resolvido na \textit{sprint} seguinte.
